Put \(N=M+1\) and \(x=N\alpha/2\). Since \(N\) is an integer, the elementary identity \(\lceil N-x\rceil=N-\lfloor x\rfloor\) gives
\[
N-k=N-\left\lceil N\left(1-\frac{\alpha}{2}\right)\right\rceil=\left\lfloor\frac{N\alpha}{2}\right\rfloor.
\]
By \(\texttt{def:finite-rank-error}\),
\[
\rho_{M,\alpha}=\frac{N-k}{N}=\frac{\lfloor N\alpha/2\rfloor}{N}\le \frac{N\alpha/2}{N}=\frac{\alpha}{2}.
\]
For an integer \(j\ge0\), the equality \(\lfloor N\alpha/2\rfloor=j\) holds if and only if
\[
j\le\frac{N\alpha}{2}<j+1,
\]
which is equivalent to \(\alpha\in[2j/N,2(j+1)/N)\). Restricting to those intervals that intersect the declared range \(0<\alpha<1\) proves the staircase assertion. Finally, \(k=M+1=N\) if and only if \(N-k=0\), equivalently \(\rho_{M,\alpha}=0\). Because \(\alpha>0\), this is equivalent to \(0<N\alpha/2<1\), hence to \(0<\alpha<2/(M+1)\).